\chapter{状态观测外力估计与模型残差}
\label{ch:observers}

编码器给出位置，却不会直接告诉我们干净的速度；电流可以换算成电机力矩，却不会直接告诉我们有人推了末端多少牛；一个新负载可能表现得很像持续外力。观测器解决的是“从已有输入输出推断看不见的量”，而不是把所有未知因素自动区分开。本章逐步推导状态观测器、卡尔曼滤波、DOB、ESO 和动量观测器，并明确它们各自估计的到底是什么。

\section{状态观测器是带测量纠正的模型}

对 $\dot x=Ax+Bu$、$y=Cx$，如果直接在计算机里运行 $\dot{\hat x}=A\hat x+Bu$，初值或模型有偏差时，估计未必能追上真实状态。加入输出残差 $y-C\hat x$：
\begin{equation}
\dot{\hat x}=A\hat x+Bu+L(y-C\hat x).
\end{equation}
$\hat x$ 表示估计，$L$ 是观测增益。定义估计误差 $\tilde x=x-\hat x$，真实方程减去估计方程：
\begin{equation}
\dot{\tilde x}=(A-LC)\tilde x.
\end{equation}
输入项抵消，因为观测器和真实系统使用了同一个输入。若观测器收到的是“期望力矩”，真实执行器却因饱和只输出一部分，输入不再相同，该差额会进入估计误差。记录实际施加的命令或可靠的实测输入因此很重要。

\begin{examplebox}[算例 14A 用位置估计关节速度]
教学关节 $A=\left[\begin{smallmatrix}0&1\\0&-2\end{smallmatrix}\right]$、$C=[1,0]$，设 $L=[l_1,l_2]^\top$，则
\[
A-LC=\begin{bmatrix}-l_1&1\\-l_2&-2\end{bmatrix}.
\]
其特征多项式为
\[
s^2+(l_1+2)s+(2l_1+l_2).
\]
希望误差极点为 $-30,-40$，期望多项式为 $s^2+70s+1200$，因此 $l_1=68$、$l_2=1200-136=1064$。

观测器两行分别为
\[
\dot{\hat q}=\hat v+68(q-\hat q),
\]
\[
\dot{\hat v}=-2\hat v+10u+1064(q-\hat q).
\]
如果位置残差是 $0.001\,\mathrm{rad}$，修正项分别为 $0.068\,\mathrm{rad/s}$ 和 $1.064\,\mathrm{rad/s^2}$。噪声也会经过同一增益放大，因此把极点设得更快并非免费。
\end{examplebox}

普通龙伯格观测器估计的是模型中已有的状态。上例只有位置和速度，没有新增“未知负载”状态，所以它不能保证独立估计并消除任意常值扰动。若需要估扰，必须增广模型或采用明确的扰动估计结构。也不能说“凡是观测器都假设扰动缓变”：缓变假设适用于某些扰动模型，不是所有状态观测问题的定义。

\subsection{分离原理的边界}

理想线性模型下，用 $u=-K\hat x$ 反馈，组合系统在适当坐标中呈块三角结构，极点由 $A-BK$ 和 $A-LC$ 两部分组成。这称为分离原理。它允许先设计反馈、再设计观测器；但不保证饱和、时延、模型失配和非线性下仍能任意独立整定。观测器“比控制器快若干倍”是常见起点，最终仍需检查噪声放大和离散极点。

\section{卡尔曼滤波把可信度放进计算}

设离散模型为
\[
x_{k+1}=F x_k+B_du_k+w_k,\qquad z_k=Hx_k+v_k.
\]
$w_k$ 是过程噪声，表示模型没有解释的状态变化；$v_k$ 是测量噪声。协方差 $Q_w=\mathbb E[w_kw_k^\top]$ 和 $R_v=\mathbb E[v_kv_k^\top]$ 描述不确定度。协方差的对角元是方差，单位是相应状态单位的平方；非对角元描述误差之间的相关性。

在常见的零均值、白噪声、不相关等假设下，预测和更新分别为
\begin{align}
\hat x_k^-&=F\hat x_{k-1}+B_du_{k-1},\\
P_k^-&=FP_{k-1}F^\top+Q_w,\\
S_k&=HP_k^-H^\top+R_v,\\
K_k&=P_k^-H^\top S_k^{-1},\\
\hat x_k&=\hat x_k^-+K_k(z_k-H\hat x_k^-).
\label{eq:cf_kalman}
\end{align}
上标 $-$ 表示还没有吸收当前测量的预测；$P$ 是估计误差协方差；$S_k$ 是新息协方差；$z_k-H\hat x_k^-$ 叫新息，即“本次测量带来的新信息”。$K_k$ 是随可信度变化的修正权重。

精确算术下可写 $P_k=(I-K_kH)P_k^-$。浮点实现常用 Joseph 形式保持对称半正定更可靠：
\begin{equation}
P_k=(I-K_kH)P_k^-(I-K_kH)^\top+K_kR_vK_k^\top.
\end{equation}
线性高斯条件下，卡尔曼滤波给出条件均值意义下的最优估计；即使噪声不是高斯，在相应二阶统计假设下仍有最佳线性估计的意义。不能把“高斯”当作能不能使用算法的简单开关，也不能把错误的噪声协方差当成最优保证。

\begin{examplebox}[算例 14B 两拍标量卡尔曼滤波]
估计一个近似不变的标量，$F=H=1$、$Q_w=0.01$、$R_v=0.25$，初值 $\hat x_0=0$、$P_0=1$。第一拍测量 $z_1=1.2$：
\[
P_1^-=1.01,\qquad K_1=\frac{1.01}{1.01+0.25}=0.801587,
\]
\[
\hat x_1=0+0.801587\times1.2=0.961905,
\qquad P_1=0.200397.
\]
第二拍测量 $z_2=0.9$：
\[
P_2^-=0.210397,\quad K_2=0.456990,
\]
\[
\hat x_2=0.961905+0.456990(0.9-0.961905)=0.933615,
\]
\[
P_2=0.114248.
\]
第一拍估计很不确定，较相信测量；吸收信息后，第二拍不再完全跟着测量走。若模型持续允许过程变化，稳态增益不会降为零。此例稳态后验方差为
\[
P=\frac{-Q_w+\sqrt{Q_w^2+4Q_wR_v}}2=0.0452494,
\]
稳态增益约为 0.180998。
\end{examplebox}

\subsection{为什么位置测量也能修正速度}

匀速模型取 $x=[q,v]^\top$，$F=\left[\begin{smallmatrix}1&T_s\\0&1\end{smallmatrix}\right]$、$H=[1,0]$。预测位置依赖速度，所以一段时间后位置误差与速度误差相关，$P$ 的非对角项不再为零。新的位置测量因此可以同时修正位置和速度。

\begin{examplebox}[算例 14C 矩阵卡尔曼的一次完整更新]
设 $T_s=0.1\,\mathrm{s}$、$\hat x=[0,1]^\top$、$P=\operatorname{diag}(0.01,0.04)$、$Q_w=\operatorname{diag}(0.0001,0.001)$、$R_v=0.0025$。这里是为了方便手算选择的教学协方差，不是编码器通用参数。

预测得到
\[
\hat x^-=[0.1,1]^\top,
\qquad P^-=\begin{bmatrix}0.0105&0.004\\0.004&0.041\end{bmatrix}.
\]
$S=0.0105+0.0025=0.013$，所以
\[
K=\frac1{0.013}\begin{bmatrix}0.0105\\0.004\end{bmatrix}
=\begin{bmatrix}0.807692\\0.307692\end{bmatrix}.
\]
测得位置 $z=0.12$，新息为 0.02。更新后
\[
\hat x=\begin{bmatrix}0.116154\\1.006154\end{bmatrix},
\quad P=\begin{bmatrix}0.00201923&0.00076923\\0.00076923&0.03976923\end{bmatrix}.
\]
测量比预测走得更远，滤波器除了增加位置估计，还小幅提高速度估计。这就是矩阵中的相关项在发挥作用。
\end{examplebox}

若把每一拍未知加速度 $a_k$ 视为方差 $\sigma_a^2$ 的随机常值，状态变化为 $[T_s^2/2,T_s]^\top a_k$，于是过程协方差为
\[
Q_w=\sigma_a^2\begin{bmatrix}T_s^4/4&T_s^3/2\\T_s^3/2&T_s^2\end{bmatrix}.
\]
若改用连续白噪声加速度驱动，积分后则会出现 $T_s^3/3$、$T_s^2/2$、$T_s$ 等系数。两式对应不同随机过程假设，不能不解释就混用。

\begin{figure}[htbp]
\centering\includegraphics[width=.95\linewidth]{figures/cf_kalman.pdf}
\caption{相同有噪位置数据的测速比较。直接差分放大噪声，低通带来滞后，卡尔曼依赖运动模型。曲线是固定随机种子的教学仿真，不表示卡尔曼对任何运动都能同时做到零噪声和零延迟。}
\end{figure}

EKF 对非线性状态方程和测量方程求局部雅可比；UKF 用代表性采样点近似传播均值和协方差；误差状态滤波在局部小误差空间工作，尤其适合旋转状态。它们的增加并不取消模型、单位、时间同步和噪声统计这些基础要求。

\section{DOB 估的是输入端总残差}

考虑 $J_r\ddot q+b\dot q=u+d$。取名义对象 $G_n$，若模型精确，$G_n^{-1}q-u=d$。用稳定低通 $Q_f(s)$ 包裹该残差：
\begin{equation}
\hat d=Q_f(s)\big[G_n^{-1}(s)q-u\big].
\end{equation}
这里扰动以正号加在输入端，因此抵消时应使用 $u=u_0-\hat d$。如果一开始把方程写成 $J_r\ddot q=u-d$，估计定义和补偿符号都要一起改变。源稿中某些公式的差异就是来自没有统一这个约定。

当 $G_n=1/(J_ns^2+b_ns)$，逆模型含有二阶微分。为了使 $Q_fG_n^{-1}$ 可实现，$Q_f$ 的相对阶数至少要补足这两阶。相对阶数是分母阶数减去分子阶数；可实现并不等于没有噪声问题，整个从位置到估计扰动的高频增益仍可能很大。

如果真实惯量与名义惯量不同，残差中还包含 $(J_n-J_r)\ddot q$；摩擦、未补偿重力和实际外力也会叠加进去。DOB 可以帮助抵消总残差，却不能仅凭这个结果宣称已经测出了“人手推力”。

\section{ESO 怎样把扰动增加成一个状态}

把单关节写成
\[
\ddot q=b_0u+f(t),
\]
其中 $b_0$ 为名义输入增益，单位是角加速度除以力矩；$f$ 把重力、摩擦、模型差异和外力等统称为角加速度形式的总扰动。增加状态 $x_3=f$，并假设其变化率在关注频段内可控制，则线性扩张状态观测器为
\begin{align}
\dot z_1&=z_2+\beta_1(q-z_1),\\
\dot z_2&=z_3+b_0u+\beta_2(q-z_1),\\
\dot z_3&=\beta_3(q-z_1).
\end{align}
忽略 $\dot f$ 时，观测误差特征式为 $s^3+\beta_1s^2+\beta_2s+\beta_3$。希望三个极点都在 $-\omega_o$，展开 $(s+\omega_o)^3$ 得
\[
\beta_1=3\omega_o,\quad\beta_2=3\omega_o^2,\quad\beta_3=\omega_o^3.
\]
因此带宽增大时，第三个增益按立方增加，这是噪声很快变严重的原因之一。

\begin{examplebox}[算例 14D ESO 补偿如何换回力矩]
取 $b_0=10$、$\omega_o=30\,\mathrm{rad/s}$，得到增益 $(90,2700,27000)$。定点目标 $q_d=0.1$，当前估计 $z_1=0.12$、$z_2=-0.05$、$z_3=1\,\mathrm{rad/s^2}$。期望误差动力学增益为 $k_p=100\,\mathrm{s^{-2}}$、$k_v=20\,\mathrm{s^{-1}}$，则
\[
u=\frac{-k_p(z_1-q_d)-k_vz_2-z_3}{b_0}
=\frac{-2+1-1}{10}=-0.2\,\mathrm{N\,m}.
\]
不能把 $z_3=1$ 直接当作 $1\,\mathrm{N\,m}$ 加到力矩上；它的单位是角加速度，必须除以 $b_0$。在 5 ms 周期下还要检查观测器离散实现，连续极点配置本身不处理实际时延和采样。
\end{examplebox}

\section{动量观测器怎样避开加速度求导}

采用全书统一动力学约定：
\begin{equation}
M(q)\ddot q+C(q,\dot q)\dot q+g(q)+\tau_f
=\tau_{\rm cmd}+\tau_{\rm ext}.
\end{equation}
$\tau_{\rm ext}$ 是环境作用于机器人的关节广义力矩。直接反解外力应为
\begin{equation}
\tau_{\rm ext}=M\ddot q+C\dot q+g+\tau_f-\tau_{\rm cmd}.
\label{eq:cf_direct_external}
\end{equation}
请注意它与“实测驱动力矩减模型所需力矩”的顺序相反。若传感器输出定义成机器人对环境的反作用力，符号会再反过来，所以必须先定义物理方向。

定义广义动量 $p=M(q)\dot q$。对时间求导：
\begin{align}
\dot p&=M\ddot q+\dot M\dot q\\
&=\tau_{\rm cmd}+\tau_{\rm ext}-g-\tau_f+(\dot M-C)\dot q.
\end{align}
对满足机器人动力学标准性质的 $C$，$\dot M-2C$ 反对称，等价地 $\dot M=C+C^\top$，因此
\begin{equation}
\dot p=\tau_{\rm cmd}+C^\top\dot q-g-\tau_f+\tau_{\rm ext}.
\end{equation}
这一步消去了显式加速度，代价是必须有质量矩阵、速度和一致的动力学模型。

定义名义动量变化项 $h_p=\tau_{\rm cmd}+\hat C^\top\dot q-\hat g-\hat\tau_f$，构造
\begin{equation}
r(t)=K_o\left[p(t)-p(0)-\int_0^t(h_p(\xi)+r(\xi))\,d\xi\right].
\end{equation}
$K_o$ 可取正对角阵，单位为 $\mathrm{s^{-1}}$。模型精确且 $K_o$ 恒定时，对上式求导：
\begin{align}
\dot r&=K_o(\dot p-h_p-r)\\
&=K_o(\tau_{\rm ext}-r).
\end{align}
所以每一轴的残差就是外力矩的一阶低通，而不是外力矩的瞬时精确测量。若初始时刻并非静止，$p(0)$ 不能擅自置零。

\begin{examplebox}[算例 14E 从观测器带宽算出检测滞后]
对一轴取 $K_o=20\,\mathrm{s^{-1}}$，时间常数 $1/K_o=0.05\,\mathrm{s}$，对应截止频率 $K_o/(2\pi)=3.183\,\mathrm{Hz}$。从零开始受到 $2\,\mathrm{N\,m}$ 阶跃外力：
\[
r(t)=2(1-e^{-20t}).
\]
50 ms 后为 $1.2642\,\mathrm{N\,m}$，150 ms 后为 $1.9004\,\mathrm{N\,m}$。若外力只持续 5 ms，峰值只到
\[
r_{\max}=2(1-e^{-0.1})=0.1903\,\mathrm{N\,m}.
\]
短冲击被严重削峰。理想一阶滤波的直流增益为 1，在完整时间区间积分、初末状态回零等条件下可以保留总面积，但这不等于峰值、接触时间或实时冲击强度都准确。
\end{examplebox}

\begin{figure}[htbp]
\centering\includegraphics[width=.95\linewidth]{figures/cf_momentum.pdf}
\caption{动量残差对阶跃和短脉冲的响应。增益提高能缩短滞后，同时也会让模型和速度测量噪声更明显。}
\end{figure}

一种实现方法是保存内部动量估计 $\hat p$，用 $r=K_o(p-\hat p)$，再积分 $\dot{\hat p}=h_p+r$。若使用前向欧拉，对理想单轴误差子系统有 $0<K_oT_s<2$ 的数值稳定要求；这只是最基本条件，不代表接近边界时具有可用性能。状态、输入和模型值还必须来自相容的时间点。

\subsection{模型误差会伪装成外力}

若模型有残差 $\Delta$，估计方程变成 $\dot r=K_o(\tau_{\rm ext}+\Delta-r)$。漏掉一个负载、摩擦方向错了、温度改变力矩常数、输出电流饱和，都可能成为 $r$ 的持续偏置。

例如没有外接触，但漏掉的重力项在某轴产生 $0.8\,\mathrm{N\,m}$，观测器最终也会报告约 $0.8\,\mathrm{N\,m}$ 的残差。静止调零能去掉当前姿态的偏置，却不能保证换姿态后仍然正确，更不能在本来有接触时把接触“调没”。

\section{关节残差怎样换成末端力}

若外力确实只作用在已知末端点，而且力旋量与雅可比在同一坐标系、同一参考点表达，则
\[
\tau_{\rm ext}=J^\top F_{\rm ext}.
\]
可用最小二乘求等效末端力
\[
\hat F=(J^\top)^+r.
\]
它回答的是“哪一个末端力最能解释这些关节残差”。如果推的是肘部、同时存在多个接触，或者残差含有零空间成分，这个结果未必是真实接触力。

\begin{examplebox}[算例 14F 两连杆末端力反算]
在标准教学位形 $q=(0,\pi/2)$，
\[
J=\begin{bmatrix}-0.5&-0.5\\0.5&0\end{bmatrix}.
\]
假设残差为 $r=[-1.5,-1]^{\top}\,\mathrm{N\,m}$。由
\[
-0.5f_x+0.5f_y=-1.5,\qquad -0.5f_x=-1
\]
先得 $f_x=2\,\mathrm{N}$，再得 $f_y=-1\,\mathrm{N}$。回代 $J^\top[2,-1]^\top=[-1.5,-1]^\top$，计算闭合。

对七轴臂应同时检查重构残差 $r-J^\top\hat F$。它明显不为零，说明“一个末端力解释全部残差”的假设不充分，而不应强行把全部残差都解释为末端接触。
\end{examplebox}

\section{从残差到事件还差哪些判断}

阈值必须基于无接触轨迹中的残差分布，并随姿态、速度、加速度和温度等因素评估。若理想独立高斯样本超过双侧 $3\sigma$ 的概率约为 0.27\%，200 Hz 下每秒期望超阈样本数约为 $200\times0.0027=0.54$。这还不是误报事件率：连续样本通常相关，去抖逻辑也会把多个样本合并成一个事件。

因此要区分偏置、瞬间尖峰和持续接触。延长连续超阈要求可以减少某类误触发，却必然增加检测时间；抬高阈值会减少误报，也会漏掉较小接触。没有一组通用阈值能绕过模型和数据质量。

\section{计算练习与逐步解答}

\textbf{练习 1}\quad 教学关节的观测器希望极点为 $-10,-20$，求 $L$。

\textbf{解}\quad 目标多项式为 $s^2+30s+200$。由 $l_1+2=30$ 得 $l_1=28$；由 $2l_1+l_2=200$ 得 $l_2=144$。比算例14A的增益小，但观测误差衰减也更慢。

\textbf{练习 2}\quad 标量先验均值为 2，方差为 0.04；测量为 2.3，测量方差为 0.09。求本次更新结果。

\textbf{解}\quad $K=0.04/(0.04+0.09)=4/13=0.307692$。后验均值为 $2+(4/13)\times0.3=2.092308$；后验方差为 $(1-4/13)0.04=0.0276923$。测量方差大于先验方差，所以更新只向测量移动约31\%。

\textbf{练习 3}\quad $K_o=40\,\mathrm{s^{-1}}$，$1\,\mathrm{N\,m}$ 阶跃外力从零开始，残差何时达到 95\%？

\textbf{解}\quad $1-e^{-40t}=0.95$，故 $t=-\ln(0.05)/40=0.074893\,\mathrm{s}$，约75 ms。把它称作40 Hz会错，角频率参数40每秒对应截止频率仅 $40/(2\pi)=6.366\,\mathrm{Hz}$。

\textbf{练习 4}\quad 单关节 $M\ddot q=0.2$、$C\dot q=0.05$、$g=1.1$、$\tau_f=0.15$、$\tau_{\rm cmd}=1.3$，单位均为力矩。求环境对机器人的外力矩。

\textbf{解}\quad 由式\eqref{eq:cf_direct_external}，$\tau_{\rm ext}=0.2+0.05+1.1+0.15-1.3=0.2\,\mathrm{N\,m}$。若反过来用驱动力矩减模型力矩，得到的是相反符号；只有明确把结果定义为机器人对环境的反作用时才能这样使用。

\section*{本章与原稿的对应}
本章整合 [A，第16--18、43、50--51页] 与 [B，第56--58、74--78页]。统一了外力符号，补足矩阵卡尔曼计算与动量观测器推导，并区分状态、总扰动、真实接触力和可解释的等效末端力。
